\subsection{Polarization and metrized operations in arbitrary degree}
\label{sec:higher-metrized-operations}

For a homogeneous polynomial $F$ of degree $d\ge3$, its scalar
polarization is the symmetric $d$-linear form
$A_F=D^dF/d!$, characterized by $A_F(x,\ldots,x)=F(x)$
\cite[Section~2.6.4, equation~(2.6.5)]{Landsberg2012}.
Here $D^dF$ is the constant $d$th differential. Its evaluation is
given by the polarization formula
\[
 D^dF[x_1,\ldots,x_d]
 =\sum_{I\subseteq\{1,\ldots,d\}}(-1)^{d-|I|}
     F\Bigl(\sum_{i\in I}x_i\Bigr).
\]
The alternating sum cancels terms missing some vector argument;
homogeneity leaves the mixed derivative of order $d$.

Using the Euclidean inner product, define the symmetric
$(d-1)$-linear map $\mu:V^{d-1}\to V$ by
\[
 \langle\mu(x_1,\ldots,x_{d-1}),x_d\rangle
       =dA_F(x_1,\ldots,x_d)
       =\frac{1}{(d-1)!}D^dF[x_1,\ldots,x_d].
\]
The complete symmetry of the right side is the metric compatibility
condition. Since $DF(x)[v]=dA_F(x,\ldots,x,v)$, this normalization
makes $\mu(x,\ldots,x)=\nabla F(x)$. For $d=3$, the product in
\eqref{eq:cubic-polarized-product} satisfies
$\langle xy,z\rangle=6A_F(x,y,z)$, so $xy=2\mu(x,y)$;
its multiplication operator is consequently twice the one used here.
Put $P(x)=\mu(x,\ldots,x)$ and
$L_xv=\mu(x,\ldots,x,v)$, with $d-2$ copies of $x$ in the latter.
Polarization and Euler's identity give
\begin{equation}\label{eq:higher-metrized-dictionary}
 F=\frac1d\langle P,x\rangle,\qquad
 \nabla F=P,\qquad \Hess F=(d-1)L_x,\qquad
 \Delta F=(d-1)\tr L_x.
\end{equation}
Every $L_x$ is self-adjoint. With $S=|P|^2$,
\begin{equation}\label{eq:higher-metrized-minimality}
 \cL(F)=(d-1)\bigl(S\tr L_x-\langle L_xP,P\rangle\bigr).
\end{equation}
In the harmonic case, $\tr L_x=0$ for all $x$, and Bochner's formula
and differentiation of $S$ give
\begin{equation}\label{eq:higher-metrized-Bochner}
 \Delta S=2(d-1)^2\tr(L_x^2),\qquad
 \Gamma(S,S)=4(d-1)^2|L_xP|^2.
\end{equation}
Indeed, $\nabla S=2(d-1)L_xP$ and
$\Delta|\nabla F|^2=2|\Hess F|^2$ when $\Delta F=0$.
Thus $\cL(F)=-qF$ is equivalent to
\[
 d(d-1)\langle L_xP,P\rangle=q\langle P,x\rangle
\]
in this harmonic normalization. The weight $q$ has degree $2d-4$
and need not be radial.

The polarization formula is an equivalent algebraic description of
the polynomial, and the contractions above express the minimal surface
equation in terms of a symmetric metrized operation. In higher degree
this can organize the identities needed to compute
$q,S,\Delta S,\Gamma(S,S)$ in
\cref{sec:harmonic-multiplier-criterion}. It becomes useful for a sign
test when additional structure, such as Clifford relations, norm
composition \cite[Chapter~1, Appendix, pp.~21--22]{Shapiro2000}, or
spectral minor identities, makes these contractions tractable.
Polarization alone supplies neither a classification nor the required
sign inequalities. For $d>3$ the operation is $(d-1)$-linear, so the
binary-algebra methods for cubics do not transfer automatically.
